\section{Introduction}
\label{sec:intro}

When a Linux kernel API changes, the drivers that use it are converted one after another, often over months and by
different developers. The earlier conversions are free, reviewed examples for the later ones: in a pilot study of
27 API migrations in Linux v5.10--v6.10, an earlier conversion of another file by the same migration (a
\emph{sibling migration}) existed in the target's git history for 96.9\% of 8{,}818 conversions~\cite{ahmed2026siblings}.
The natural way to reuse such examples in the kernel is a Coccinelle semantic patch~\cite{padioleau2008documenting}:
a rule that kernel maintainers can read, review and apply tree-wide, unlike a stream of per-file edits from a
language model.

Writing such rules automatically from examples is an old problem. Spinfer~\cite{serrano2020spinfer} infers them
with dedicated heuristics, and recent work prompts frontier LLMs to write them~\cite{yusuf2025howfar}, finding
that LLM-written patches are often syntactically invalid and that both approaches struggle when the examples
contain several variants of a change. Frontier APIs are also a poor fit for kernel workflows: maintainers cannot
inspect them, and they have likely seen the very migrations they are evaluated on.

We ask whether \emph{small open models running on a laptop} can write correct migration rules when the work is
structured for them. We adapt guided code generation (GCG)~\cite{almorsi2025gcg}, a decomposition framework
previously shown to help an 8B model on HumanEval, to semantic patches: the training siblings are decomposed into
a tree of change variants; the model writes one small rule per variant under test feedback from the siblings
themselves; deterministic components repair mechanical SmPL errors and generalise driver-specific names; rules are
composed greedily; and a compile critic either repairs or rejects a rule whose output does not build.

On \testMigrations{} migrations whose APIs postdate all three models, the framework takes Qwen2.5-Coder 7B from
\testQwenCoderOneShot{} to \testQwenCoderFull{} solved conversions out of \testBuildable{}, close to but below
Spinfer (\testSpinfer), with the highest precision of all systems on the test set
(\testSafeQwenCoderPrecision). Decomposition and guidance carry most of the gain; the compile critic adds no
solved conversions but removes every output that matches the human fix and fails to build. The two general-purpose
models, Llama~3.1 8B and Gemma~3 4B, remain weak, mainly because they rarely write SmPL that parses. Changes that
span a declaration and its uses, such as a callback whose return type changes, defeat every local model.

\paragraph{Contributions.}
\begin{itemize}
  \item A guided framework for semantic-patch generation with 4--8B local models (Section~\ref{sec:method}), whose
    every step except rule writing is deterministic and auditable.
  \item A pre-registered test set beyond the models' training cutoffs: \testMigrations{} Linux API migrations whose new API first
    appeared after all three models' training cutoffs, with \testTargets{} held-out driver conversions
    (Section~\ref{sec:setup}).
  \item An evaluation that compiles candidate outputs on x86\_64 \texttt{allmodconfig} and reports an automatic,
    judgement-free success metric and precision, three seeds per configuration, an ablation of every component,
    and an analysis of the failures (Section~\ref{sec:results}).
\end{itemize}
All code, prompts, model responses and results are released.\footnote{\url{https://github.com/MohannedAhmed67/kernel-sibling-migrations}}
